This review compares multi-agent diagnostics and selected serving infrastructure mechanisms through a common set of operational dimensions: evidence inputs, training and replay requirements, granularity, intervention time, state compatibility, and resource cost. Their interaction cannot be summarized by a universal conflict between diagnostics and KV-cache eviction. Information is lost only when a needed, non-reconstructable diagnostic dependency is actually discarded; otherwise reclaiming state may increase recomputation or latency without degrading the retained trace.

The proposed conditional retention framework and evaluation protocol are synthesis contributions, not an already validated serving architecture. The available literature establishes the importance of trace completeness and describes serving policies that manage related, but not identical, state. It does not yet establish a general causal relationship between cache decisions and diagnostic accuracy. A rigorous next step is full-text source verification followed by matched, prospective experiments measuring attribution quality, false alarms, recovery, and serving cost under controlled retention budgets. These conditions define when reliability-aware serving may be beneficial and, equally importantly, when the proposed cross-layer trade-off is absent.